\documentclass[11pt]{article}

\usepackage[a4paper,margin=1in]{geometry}
\usepackage{amsmath,amssymb,mathtools}
\usepackage{bm}
\usepackage{enumitem}
\usepackage{hyperref}

\newcommand{\dd}{\,\mathrm d}
\newcommand{\R}{\mathbb R}
\newcommand{\OmegaD}{\Omega}
\newcommand{\V}{V}
\newcommand{\VT}{V_T}
\newcommand{\sig}{\sigma}
\newcommand{\W}{W}
\newcommand{\cvec}{\bm c}
\newcommand{\astar}{\alpha_1}
\newcommand{\bstar}{\alpha_2}
\newcommand{\phistar}{\phi^\star}
\newcommand{\phit}{\phi_T}
\newcommand{\rhot}{\rho_T}
\newcommand{\epsT}{\varepsilon_T}
\newcommand{\epsp}{\varepsilon_\Phi}
\newcommand{\etaT}{\eta_T}
\newcommand{\etap}{\eta_\Phi}
\newcommand{\Honez}{H^1_0(\OmegaD)}

\title{Closed-Loop Window Optimization for the Torsion-Initialized Newton Method}
\author{}
\date{2026-07-13}

\begin{document}
\maketitle

\section{Purpose}

This note formulates a closed-loop version of the torsion-initialized Newton
band-selection problem.  The input torsion thresholds are fixed through two parameters
\(\alpha_1\) and \(\alpha_2\).  The output thresholds \(c_1\) and \(c_2\)
are not chosen by matching a window applied to an already computed potential.
Instead, they are chosen by solving the nonlinear state equation
\[
  -\Delta \phistar = \W(\phistar;c_1,c_2,\epsp)
  \quad \text{in } \OmegaD,
  \qquad
  \phistar = 0
  \quad \text{on } \partial\OmegaD,
\]
and minimizing the \(L^2(\OmegaD)\) distance between \(\phistar\) and a
torsion-derived target potential \(\phit\).  This is the abstraction suggested
by the numerical experiments: \(c_1\) and \(c_2\) should be selected through
the semilinear fixed point they produce, not only through a pointwise fit of
the right-hand side.

\section{Notation}

The computational domain is a bounded Lipschitz set
\(\OmegaD\subset\R^2\).  Its boundary is denoted by \(\partial\OmegaD\).
All scalar potentials in this note satisfy homogeneous Dirichlet boundary
conditions.  The energy space is
\[
  \V = \Honez
  =
  \{v\in H^1(\OmegaD): v=0 \text{ on } \partial\OmegaD\}.
\]
The \(L^2(\OmegaD)\) inner product and norm are
\[
  (u,v)_{L^2(\OmegaD)}
  =
  \int_{\OmegaD} u(x)v(x)\dd x,
  \qquad
  \|u\|_{L^2(\OmegaD)}
  =
  \left((u,u)_{L^2(\OmegaD)}\right)^{1/2}.
\]

The fixed torsion parameters are
\[
  0<\alpha_1<\alpha_2<1.
\]
The torsion function \(T\in \V\) is defined by
\[
  -\Delta T = 1
  \quad \text{in } \OmegaD,
  \qquad
  T=0
  \quad \text{on } \partial\OmegaD.
\]
Equivalently, \(T\in\V\) satisfies
\[
  \int_{\OmegaD}\nabla T\cdot\nabla v\dd x
  =
  \int_{\OmegaD} v\dd x
  \qquad
  \forall v\in\V.
\]
Let
\[
  T_{\max} = \operatorname*{ess\,sup}_{x\in\OmegaD} T(x).
\]
The torsion thresholds are
\[
  c_{T,1} = \alpha_1 T_{\max},
  \qquad
  c_{T,2} = \alpha_2 T_{\max}.
\]
The torsion smoothing scale is
\[
  \epsT = \eta_T(c_{T,2}-c_{T,1}),
\]
where \(\eta_T>0\) is a prescribed dimensionless smoothing ratio.

The amplitude of the density window is denoted by \(A>0\).  In the current
scripts this amplitude is \(\texttt{rho\_amp}\), usually \(A=1\).  For every
\(\varepsilon>0\), define the clipped logistic idealization
\[
  \sig_\varepsilon(z)
  =
  \frac{1}{1+\exp(-z/\varepsilon)}.
\]
The implementation clips the argument \(z/\varepsilon\) for numerical
stability, but the mathematical model in this note uses the smooth expression
above.  Given two thresholds \(a<b\), define the two-sided window
\[
  \boxed{
  \W(s;a,b,\varepsilon)
  =
  A\left(\sig_\varepsilon(s-a)-\sig_\varepsilon(s-b)\right).
  }
\]
This window is near zero for \(s\ll a\), near \(A\) for \(a\ll s\ll b\), and
near zero again for \(s\gg b\).

\section{Torsion Target}

The torsion-designed density is
\[
  \rhot(x)
  =
  \W(T(x);c_{T,1},c_{T,2},\epsT).
\]
The target potential \(\phit\in\V\) is the Poisson response to this density:
\[
  -\Delta \phit = \rhot
  \quad \text{in } \OmegaD,
  \qquad
  \phit=0
  \quad \text{on } \partial\OmegaD.
\]
Equivalently,
\[
  \int_{\OmegaD}\nabla \phit\cdot\nabla v\dd x
  =
  \int_{\OmegaD}\rhot v\dd x
  \qquad
  \forall v\in\V.
\]
The function \(\phit\) is fixed once \(\OmegaD\), \(A\), \(\eta_T\),
\(\alpha_1\), and \(\alpha_2\) have been chosen.  It is the quantity that the
closed-loop semilinear equilibrium will try to match.

\section{Closed-Loop State Equation}

The optimization variables are the semilinear thresholds
\[
  \cvec = (c_1,c_2)\in\R^2.
\]
They live in the \(\phi\)-scale, not the torsion scale.  The admissible set is
\[
  \mathcal C
  =
  \left\{
    (c_1,c_2)\in\R^2:
    0\le c_1,\quad
    c_1+\delta_{\min}\le c_2,\quad
    c_2\le C_{\max}
  \right\}.
\]
Here \(\delta_{\min}>0\) is a prescribed minimum width, and \(C_{\max}>0\) is
a prescribed upper search bound.  A natural choice is
\[
  C_{\max} = \operatorname*{ess\,sup}_{x\in\OmegaD}\phit(x),
\]
possibly enlarged by a safety factor if the semilinear solution can exceed
\(\phit\).

For a fixed continuation stage, the semilinear smoothing scale is
\[
  \epsp(\cvec) = \eta_\Phi(c_2-c_1),
\]
where \(\eta_\Phi>0\) is the chosen continuation ratio.  In the scripts this
ratio is one of the values in the sequence
\[
  \eta_\Phi\in \{0.11,\ 0.08,\ 0.06\}.
\]
For each \(\cvec\in\mathcal C\), define the state
\(\phi(\cvec)\in\V\) as a solution of
\[
  -\Delta\phi(\cvec)
  =
  \W\!\left(
    \phi(\cvec);c_1,c_2,\epsp(\cvec)
  \right)
  \quad \text{in } \OmegaD,
  \qquad
  \phi(\cvec)=0
  \quad \text{on } \partial\OmegaD.
\]
The corresponding weak form is
\[
  \int_{\OmegaD}\nabla\phi(\cvec)\cdot\nabla v\dd x
  =
  \int_{\OmegaD}
    \W\!\left(\phi(\cvec);c_1,c_2,\epsp(\cvec)\right)v
  \dd x
  \qquad
  \forall v\in\V.
\]

The semilinear problem can have multiple solutions if the window derivative is
large enough.  Therefore the state map
\[
  S:\mathcal C\to\V,
  \qquad
  S(\cvec)=\phi(\cvec),
\]
must be understood as the branch selected by the chosen numerical procedure:
the continuation ratios, the Newton initialization, the damping strategy, and
the linear solver tolerances.  The optimization problem below is therefore a
closed-loop optimization over this selected solution branch.

\section{Reduced Optimization Problem}

The closed-loop threshold-selection problem is
\[
  \boxed{
  \min_{\cvec\in\mathcal C}
  J(\cvec)
  =
  \frac{1}{2}
  \left\|S(\cvec)-\phit\right\|_{L^2(\OmegaD)}^2 .
  }
\]
The relative version, useful for reporting, is
\[
  J_{\mathrm{rel}}(\cvec)
  =
  \frac{
    \left\|S(\cvec)-\phit\right\|_{L^2(\OmegaD)}
  }{
    \left\|\phit\right\|_{L^2(\OmegaD)}
  },
\]
provided \(\|\phit\|_{L^2(\OmegaD)}>0\).  The minimizer is denoted by
\[
  \cvec^\star = (c_1^\star,c_2^\star).
\]
The corresponding state is
\[
  \phistar = S(\cvec^\star).
\]
By construction, \(\phistar\) satisfies
\[
  -\Delta\phistar
  =
  \W(\phistar;c_1^\star,c_2^\star,\epsp(\cvec^\star)),
  \qquad
  \phistar|_{\partial\OmegaD}=0,
\]
and it is the closest semilinear equilibrium to \(\phit\) among the
admissible thresholds in the selected solution branch.

\section{First-Order Sensitivity}

This section derives the gradient of \(J\) with respect to \(c_1\) and
\(c_2\).  The derivation is useful even if the first implementation uses a
derivative-free search, because it states exactly which closed-loop quantity is
being optimized.

Let
\[
  s\in\R,\qquad
  \varepsilon>0,\qquad
  a<b.
\]
Define
\[
  \ell_a = \sig_\varepsilon(s-a),
  \qquad
  \ell_b = \sig_\varepsilon(s-b).
\]
The derivatives of the window with respect to its scalar argument and
thresholds are
\[
  \W_s(s;a,b,\varepsilon)
  =
  \frac{A}{\varepsilon}
  \left(\ell_a(1-\ell_a)-\ell_b(1-\ell_b)\right),
\]
\[
  \W_a(s;a,b,\varepsilon)
  =
  -\frac{A}{\varepsilon}\ell_a(1-\ell_a),
  \qquad
  \W_b(s;a,b,\varepsilon)
  =
  \frac{A}{\varepsilon}\ell_b(1-\ell_b).
\]
The derivative with respect to \(\varepsilon\) is
\[
  \W_\varepsilon(s;a,b,\varepsilon)
  =
  A\left[
    -\frac{s-a}{\varepsilon^2}\ell_a(1-\ell_a)
    +
    \frac{s-b}{\varepsilon^2}\ell_b(1-\ell_b)
  \right].
\]
Since
\[
  \epsp(\cvec)=\eta_\Phi(c_2-c_1),
\]
the total derivatives of the right-hand side with respect to the optimization
variables are
\[
  G_1(s;\cvec)
  =
  \W_{c_1}(s;c_1,c_2,\epsp)
  -
  \eta_\Phi\W_\varepsilon(s;c_1,c_2,\epsp),
\]
\[
  G_2(s;\cvec)
  =
  \W_{c_2}(s;c_1,c_2,\epsp)
  +
  \eta_\Phi\W_\varepsilon(s;c_1,c_2,\epsp).
\]
Here \(\W_{c_1}\) is \(\W_a\) evaluated at \(a=c_1\), and
\(\W_{c_2}\) is \(\W_b\) evaluated at \(b=c_2\).

Assume that \(\phi=S(\cvec)\) is a differentiable point of the selected state
branch.  Let
\[
  z_i = \frac{\partial S}{\partial c_i}(\cvec)\in\V,
  \qquad i\in\{1,2\}.
\]
Differentiating the state equation gives the sensitivity equations
\[
  -\Delta z_i
  -
  \W_s(\phi;c_1,c_2,\epsp) z_i
  =
  G_i(\phi;\cvec)
  \quad \text{in } \OmegaD,
  \qquad
  z_i=0
  \quad \text{on } \partial\OmegaD.
\]
In weak form,
\[
  \int_{\OmegaD}\nabla z_i\cdot\nabla v\dd x
  -
  \int_{\OmegaD}
    \W_s(\phi;c_1,c_2,\epsp) z_i v
  \dd x
  =
  \int_{\OmegaD}G_i(\phi;\cvec)v\dd x
  \qquad
  \forall v\in\V.
\]
The direct reduced gradient is
\[
  \frac{\partial J}{\partial c_i}(\cvec)
  =
  \int_{\OmegaD}(\phi-\phit)z_i\dd x,
  \qquad i\in\{1,2\}.
\]

\section{Adjoint Gradient}

Solving two sensitivity equations is cheap in this two-parameter problem, but
the adjoint form is cleaner and extends to richer parameterizations.  Define
the linearized state operator
\[
  \mathcal L_\phi z
  =
  -\Delta z
  -
  \W_s(\phi;c_1,c_2,\epsp)z,
  \qquad z\in\V.
\]
With homogeneous Dirichlet boundary conditions, this operator is formally
self-adjoint in \(L^2(\OmegaD)\).  The adjoint variable \(p\in\V\) is defined
by
\[
  \mathcal L_\phi p = \phi-\phit
  \quad \text{in } \OmegaD,
  \qquad
  p=0
  \quad \text{on } \partial\OmegaD.
\]
Equivalently,
\[
  \int_{\OmegaD}\nabla p\cdot\nabla v\dd x
  -
  \int_{\OmegaD}
    \W_s(\phi;c_1,c_2,\epsp)p v
  \dd x
  =
  \int_{\OmegaD}(\phi-\phit)v\dd x
  \qquad
  \forall v\in\V.
\]
Using the sensitivity equations,
\[
  \boxed{
  \frac{\partial J}{\partial c_i}(\cvec)
  =
  \int_{\OmegaD} p(x)G_i(\phi(x);\cvec)\dd x,
  \qquad i\in\{1,2\}.
  }
\]

For an unconstrained interior minimizer, the necessary first-order conditions
are
\[
  \frac{\partial J}{\partial c_1}(\cvec^\star)=0,
  \qquad
  \frac{\partial J}{\partial c_2}(\cvec^\star)=0.
\]
For the constrained admissible set \(\mathcal C\), the correct optimality
condition is the Karush--Kuhn--Tucker condition for the inequalities
\[
  -c_1\le 0,\qquad
  c_1+\delta_{\min}-c_2\le 0,\qquad
  c_2-C_{\max}\le 0.
\]
If none of these inequalities is active, the unconstrained equations above
hold.  If one or more inequalities are active, the gradient need not vanish;
it must be balanced by nonnegative multipliers attached to the active
constraints.

\section{Discrete Version}

Let \(\mathcal T_h\) be a conforming triangular mesh of \(\OmegaD\), and let
\[
  V_h
  =
  \{v_h\in C^0(\overline\OmegaD):
    v_h|_K\in\mathbb P_p(K)\ \forall K\in\mathcal T_h,\ 
    v_h|_{\partial\OmegaD}=0\}
\]
be the continuous Lagrange space of polynomial degree \(p\).  This is the
space used by the DOLFINx comparison runner.  The discrete torsion
\(T_h\in V_h\), torsion target density \(\rho_{T,h}\), and target potential
\(\phi_{T,h}\in V_h\) are defined by replacing \(T\), \(\rhot\), and \(\phit\)
with their finite element counterparts.

For each \(\cvec\in\mathcal C\), the discrete state
\(\phi_h(\cvec)\in V_h\) satisfies
\[
  \int_{\OmegaD}\nabla\phi_h(\cvec)\cdot\nabla v_h\dd x
  =
  \int_{\OmegaD}
    \W(\phi_h(\cvec);c_1,c_2,\epsp(\cvec))v_h
  \dd x
  \qquad
  \forall v_h\in V_h.
\]
The discrete reduced objective is
\[
  J_h(\cvec)
  =
  \frac{1}{2}
  \|\phi_h(\cvec)-\phi_{T,h}\|_{L^2(\OmegaD)}^2.
\]
All integrals in \(J_h\), the state residual, the sensitivity equations, and
the adjoint equation should be evaluated using a quadrature rule of sufficient
degree for the selected polynomial order and the steepness of
\(\W\).  Since \(\W\) is not a polynomial, increasing polynomial degree does
not by itself remove the need for adequate quadrature near the two transition
layers \(s\approx c_1\) and \(s\approx c_2\).

\section{Relation to the Current v2 Fitted Window}

The current v2 runner fits \(c_1\) and \(c_2\) with the surrogate objective
\[
  \min_{(c_1,c_2)\in\mathcal C}
  \left\|
    \W(\phi_{T,h};c_1,c_2,\epsp(c_1,c_2))
    -
    \rho_{T,h}
  \right\|_{L^2(\OmegaD)}^2.
\]
This is a pointwise right-hand-side fit on quadrature samples.  It is fast
because it does not solve the semilinear state equation for every candidate
\((c_1,c_2)\).  However, it is not the closed-loop objective above.  It can
fail for three distinct reasons:
\begin{enumerate}[label=(\roman*)]
  \item The level sets of \(\phi_{T,h}\) need not align well enough with the
  level sets of \(T_h\), even though \(\phi_{T,h}\) is generated from the
  torsion-designed density.
  \item The final semilinear state \(\phi_h(\cvec)\) need not stay close to
  \(\phi_{T,h}\) after solving the nonlinear equation
  \[
    -\Delta\phi_h=\W(\phi_h;c_1,c_2,\epsp).
  \]
  \item A right-hand-side \(L^2\) fit does not enforce mass, active-set
  overlap, plateau-set overlap, or closeness of the final potential.
\end{enumerate}
The closed-loop objective removes the second mismatch by evaluating the
candidate thresholds through the nonlinear PDE itself.

\section{Possible Algorithms}

\subsection{Coarse Closed-Loop Search}

The most robust first implementation is a coarse reduced search:
\begin{enumerate}[label=\arabic*.]
  \item Fix \(\alpha_1\), \(\alpha_2\), \(\eta_T\), and \(\eta_\Phi\).
  \item Compute \(T_h\), \(\rho_{T,h}\), and \(\phi_{T,h}\).
  \item Choose a coarse grid of admissible pairs \((c_1,c_2)\in\mathcal C\).
  \item For each pair, solve the semilinear state equation to obtain
  \(\phi_h(c_1,c_2)\).
  \item Score each pair with \(J_h(c_1,c_2)\).
  \item Refine the search around the best pair.
\end{enumerate}
This method is expensive but directly measures the desired quantity.  It is
appropriate for validating the formulation and producing reference data.

\subsection{Adjoint Reduced Optimization}

Once the objective is validated, one can use the adjoint gradient:
\begin{enumerate}[label=\arabic*.]
  \item For the current pair \(\cvec\), solve the state equation for
  \(\phi_h(\cvec)\).
  \item Solve the adjoint equation for \(p_h\).
  \item Assemble the gradient components
  \[
    \frac{\partial J_h}{\partial c_i}
    =
    \int_{\OmegaD}p_hG_i(\phi_h;\cvec)\dd x,
    \qquad i=1,2.
  \]
  \item Take a constrained optimization step in \(\mathcal C\).
  \item Repeat until the projected gradient and objective decrease satisfy
  prescribed tolerances.
\end{enumerate}
This is the natural method if the same idea is later extended from two
thresholds to a larger parameterization of the nonlinear window.

\subsection{Fixed Smoothing First, Continuation Later}

The optimization problem should first be posed at one fixed value of
\(\eta_\Phi\).  That is, during one outer optimization run we fix
\[
  \epsp(\cvec)=\eta_\Phi(c_2-c_1)
\]
with a single chosen value of \(\eta_\Phi\), for example
\(\eta_\Phi=0.06\).  The outer optimizer then sees a stationary objective
\[
  J_{h,\eta_\Phi}(\cvec)
  =
  \frac{1}{2}
  \|\phi_{h,\eta_\Phi}(\cvec)-\phi_{T,h}\|_{L^2(\OmegaD)}^2,
\]
where \(\phi_{h,\eta_\Phi}(\cvec)\) solves the semilinear problem with that
fixed smoothing ratio.

Only after this fixed-\(\eta_\Phi\) problem is understood should one run a
continuation sequence.  A continuation sequence is a family of optimization
problems
\[
  \eta_\Phi^{(0)}>\eta_\Phi^{(1)}>\cdots>\eta_\Phi^{(m)}
\]
where the optimizer for stage \(j+1\) is initialized with the best pair found
at stage \(j\).  This is a continuation in the outer reduced problem, not only
in the inner Newton solve.  It avoids mixing two effects in the first test:
\begin{enumerate}[label=(\roman*)]
  \item whether the closed-loop objective identifies useful thresholds, and
  \item whether the sharp-window continuation path is numerically stable.
\end{enumerate}

\subsection{Inexact Reduced Optimization}

The exact reduced method assumes that, for every candidate
\(\cvec=(c_1,c_2)\), the state equation is solved exactly.  In computation,
this is neither possible nor necessary.  Let the discrete semilinear residual
be the functional
\[
  F_h(\phi_h,\cvec;v_h)
  =
  \int_{\OmegaD}\nabla\phi_h\cdot\nabla v_h\dd x
  -
  \int_{\OmegaD}
    \W(\phi_h;c_1,c_2,\epsp(\cvec))v_h
  \dd x,
  \qquad v_h\in V_h.
\]
The exact state equation is
\[
  F_h(\phi_h,\cvec;v_h)=0
  \qquad
  \forall v_h\in V_h.
\]
An inexact reduced method replaces the exact state
\(\phi_h(\cvec)\) by an approximate state \(\widetilde\phi_h(\cvec)\) such
that
\[
  \|F_h(\widetilde\phi_h,\cvec;\cdot)\|_{V_h'} \le \tau_{\mathrm{state}},
\]
where \(V_h'\) is the algebraic dual of \(V_h\), and
\(\tau_{\mathrm{state}}\) is a chosen state tolerance.  The norm can be the
energy-dual norm
\[
  \|F_h\|_{V_h'}^2
  =
  \bm r^T \bm A^{-1}\bm r,
\]
where \(\bm r\) is the residual vector and \(\bm A\) is the stiffness matrix,
or a cheaper Euclidean residual norm for early experiments.

The important point is that the inner Newton solve does not need to reach
machine precision.  It only needs to be accurate enough that the error in the
reduced objective and adjoint gradient is smaller than the optimization
signal.  A practical rule is
\[
  \|F_h(\widetilde\phi_h,\cvec;\cdot)\|_{V_h'}
  \le
  \theta_k
  \max\{1,\|\nabla_{\cvec}J_h(\cvec)\|_2\},
\]
where \(0<\theta_k<1\) is tightened as the outer optimization converges.
Early outer iterations can use loose Newton tolerances.  Near a candidate
minimum, the state tolerance must be tightened; otherwise the adjoint gradient
can point in a direction dominated by state-solve error.

Operationally, this suggests the following first implementation:
\begin{enumerate}[label=\arabic*.]
  \item Fix one \(\eta_\Phi\).
  \item Start from the v2 fitted pair or from a successful band-study pair.
  \item For each outer step, warm-start Newton from the previous semilinear
  state.
  \item Stop the inner Newton solve when the state residual is below an
  outer-iteration-dependent tolerance, not necessarily machine precision.
  \item Compute the adjoint and gradient at the approximate state.
  \item Take a line-searched constrained step in \(\mathcal C\).
  \item If the line search becomes ambiguous, recompute the current state with
  a stricter tolerance and retry.
\end{enumerate}

\subsection{All-at-Once Formulation in \texorpdfstring{\(\mathcal C\times V_h\)}{C x Vh}}

It is natural to ask whether one can avoid fully solving the Newton problem
for each \(\cvec\) by optimizing directly over both the thresholds and the
state:
\[
  (\cvec,\phi_h)\in\mathcal C\times V_h.
\]
The mathematically faithful all-at-once problem is the equality-constrained
optimization problem
\[
  \boxed{
  \begin{aligned}
  \min_{\cvec\in\mathcal C,\ \phi_h\in V_h}
  \quad&
  \frac{1}{2}\|\phi_h-\phi_{T,h}\|_{L^2(\OmegaD)}^2,\\
  \text{subject to}\quad&
  F_h(\phi_h,\cvec;v_h)=0
  \qquad \forall v_h\in V_h.
  \end{aligned}
  }
\]
This formulation does not require an inner nonlinear solve at every outer
threshold update.  Instead, one solves the KKT system for
\((\phi_h,\cvec,p_h)\), where \(p_h\in V_h\) is the adjoint multiplier for the
state equation.  The first-order equations are:
\[
  F_h(\phi_h,\cvec;v_h)=0
  \qquad
  \forall v_h\in V_h,
\]
\[
  \int_{\OmegaD}\nabla p_h\cdot\nabla w_h\dd x
  -
  \int_{\OmegaD}
    \W_s(\phi_h;c_1,c_2,\epsp)p_h w_h
  \dd x
  =
  \int_{\OmegaD}(\phi_h-\phi_{T,h})w_h\dd x
  \qquad
  \forall w_h\in V_h,
\]
and
\[
  \int_{\OmegaD}p_hG_i(\phi_h;\cvec)\dd x
  +
  \text{constraint terms for }\mathcal C
  =
  0,
  \qquad i=1,2.
\]
This is equivalent to the reduced adjoint optimality system when the state
equation is exactly satisfied.  A Newton or SQP method applied to this KKT
system would update \(\phi_h\), \(p_h\), and \(\cvec\) simultaneously.

There is also a simpler penalty alternative:
\[
  \min_{\cvec\in\mathcal C,\ \phi_h\in V_h}
  \frac{1}{2}\|\phi_h-\phi_{T,h}\|_{L^2(\OmegaD)}^2
  +
  \frac{\gamma}{2}
  \|F_h(\phi_h,\cvec;\cdot)\|_{V_h'}^2.
\]
For finite \(\gamma\), this is not the same problem as the constrained
closed-loop formulation.  A minimizer may trade state-equation residual against
potential matching.  This can be useful as a globalization or diagnostic
device, but it should not be interpreted as the final torsion-initialized Newton threshold
selection problem unless the residual term is driven sufficiently close to
zero.

The all-at-once approach is attractive because it may avoid oversolving the
state equation for poor threshold pairs.  Its cost is complexity: the KKT
matrix is larger, generally indefinite, and its globalization is more delicate
than a line-searched reduced method.  For this reason, the recommended first
implementation is an inexact reduced adjoint method.  If that method shows the
right trend but spends too much time in inner Newton solves, then an all-at-once
SQP or residual-penalty method becomes the next reasonable abstraction.

\subsection{Residual-Penalty All-at-Once Prototype}

The first all-at-once implementation should use a residual-penalty objective
rather than the full KKT system.  This is not the final constrained problem,
but it is much simpler to debug and it tests whether simultaneous updates in
\(\mathcal C\times V_h\) can find useful window thresholds without repeatedly
solving the semilinear state equation to high accuracy.

Use center-width variables
\[
  m=\frac{c_1+c_2}{2},
  \qquad
  w=c_2-c_1,
\]
so that
\[
  c_1=m-\frac{w}{2},
  \qquad
  c_2=m+\frac{w}{2},
  \qquad
  \epsp=\eta_\Phi w.
\]
The admissible set is enforced by projection:
\[
  w_{\min}\le w\le C_{\max},
  \qquad
  \frac{w}{2}\le m\le C_{\max}-\frac{w}{2}.
\]

Let \(\bm r(\phi_h,m,w)\) be the assembled residual vector associated with
\[
  F_h(\phi_h,m,w;v_h)
  =
  \int_{\OmegaD}\nabla\phi_h\cdot\nabla v_h\dd x
  -
  \int_{\OmegaD}\W(\phi_h;c_1,c_2,\eta_\Phi w)v_h\dd x.
\]
The prototype minimizes
\[
  \boxed{
  \begin{aligned}
  \mathcal J_\gamma(\phi_h,m,w)
  &=
  \frac{\lambda_\phi}{2}
  \|\phi_h-\phi_{T,h}\|_{L^2(\OmegaD)}^2\\
  &\quad+
  \frac{\lambda_\rho}{2}
  \|\rho_h-\rho_{T,h}\|_{L^2(\OmegaD)}^2\\
  &\quad+
  \frac{\lambda_M}{2}
  \left(\int_{\OmegaD}(\rho_h-\rho_{T,h})\dd x\right)^2\\
  &\quad+
  \frac{\gamma}{2}\|\bm r(\phi_h,m,w)\|_2^2,
  \end{aligned}
  }
\]
where
\[
  \rho_h=\W(\phi_h;c_1,c_2,\eta_\Phi w).
\]
The weights satisfy
\[
  \lambda_\phi\ge0,\qquad
  \lambda_\rho\ge0,\qquad
  \lambda_M\ge0,\qquad
  \gamma>0.
\]
The density and mass terms are included because the intended practical target
is not only closeness of potentials, but also closeness of the designed band
and the final semilinear band.  Active-set and plateau-set areas should remain
diagnostics in this first prototype because their indicator definitions are
nonsmooth.

The scalar derivatives needed for the \(m,w\) update are
\[
  \W_m
  =
  \W_{c_1}+\W_{c_2},
\]
and
\[
  \W_w
  =
  -\frac{1}{2}\W_{c_1}
  +
  \frac{1}{2}\W_{c_2}
  +
  \eta_\Phi\W_\varepsilon.
\]
The residual-penalty derivatives with respect to \(m\) and \(w\) are
\[
  \frac{\partial}{\partial m}
  \frac{\gamma}{2}\|\bm r\|_2^2
  =
  \gamma\,\bm r^T\bm r_m,
  \qquad
  \frac{\partial}{\partial w}
  \frac{\gamma}{2}\|\bm r\|_2^2
  =
  \gamma\,\bm r^T\bm r_w,
\]
where \(\bm r_m\) and \(\bm r_w\) are the assembled vectors for
\[
  -\int_{\OmegaD}\W_m(\phi_h;m,w)v_h\dd x,
  \qquad
  -\int_{\OmegaD}\W_w(\phi_h;m,w)v_h\dd x.
\]
The coefficient gradient with respect to \(\phi_h\) is
\[
  \bm g_\phi
  =
  \bm g_{\mathrm{data}}
  +
  \gamma \bm J_\phi^T\bm r,
\]
where \(\bm J_\phi\) is the Jacobian of \(\bm r\) with respect to
\(\phi_h\).  In weak form,
\[
  \bm J_\phi
  \leftrightarrow
  \int_{\OmegaD}\nabla z_h\cdot\nabla v_h\dd x
  -
  \int_{\OmegaD}\W_s(\phi_h;c_1,c_2,\eta_\Phi w)z_hv_h\dd x.
\]

Use a preconditioned descent direction for the state variable:
\[
  \bm P\bm d_\phi=-\bm g_\phi,
\]
with
\[
  \bm P
  \leftrightarrow
  \int_{\OmegaD}\nabla z_h\cdot\nabla v_h\dd x
  +
  \mu_P\int_{\OmegaD}z_hv_h\dd x.
\]
The threshold direction is the negative scalar gradient in \((m,w)\), scaled
so that the un-line-searched threshold update is not larger than a prescribed
fraction of \(C_{\max}\).  A backtracking Armijo line search is then applied
to the combined step
\[
  (\phi_h,m,w)
  \leftarrow
  (\phi_h,m,w)
  +
  \alpha(d_\phi,d_m,d_w),
\]
with projection of \((m,w)\) after each trial update.

Every iteration should report both objective quality and state feasibility:
\[
  \mathcal J_\gamma,\quad
  \|\phi_h-\phi_{T,h}\|_{L^2},\quad
  \|\rho_h-\rho_{T,h}\|_{L^2},\quad
  \left|\int_{\OmegaD}(\rho_h-\rho_{T,h})\dd x\right|,
  \quad
  \|\bm r\|_2.
\]
If the band metrics improve but \(\|\bm r\|_2\) remains large, the iterate is
not a valid semilinear equilibrium.  It is only a useful search state.

\subsection{Composite Diagnostics}

The pure potential objective is
\(\|\phi_h-\phi_{T,h}\|_{L^2(\OmegaD)}\).  If the scientific goal is closer to
matching density geometry, then one can add terms:
\[
  J_{\mathrm{comp}}(\cvec)
  =
  \frac{1}{2}\|\phi_h(\cvec)-\phi_{T,h}\|_{L^2(\OmegaD)}^2
  +
  \frac{\lambda_m}{2}
  \left(\int_{\OmegaD}\rho_h(\cvec)-\rho_{T,h}\dd x\right)^2
  +
  \frac{\lambda_\rho}{2}
  \|\rho_h(\cvec)-\rho_{T,h}\|_{L^2(\OmegaD)}^2.
\]
Here
\[
  \rho_h(\cvec)
  =
  \W(\phi_h(\cvec);c_1,c_2,\epsp(\cvec)),
\]
and \(\lambda_m\ge 0\), \(\lambda_\rho\ge 0\) are user-chosen weights.  Active
set and plateau set overlaps are useful diagnostics, but their sharp
indicator definitions are nonsmooth.  If they are used in optimization rather
than only reporting, they should be replaced by smooth approximations.

\section{Implementation Implications}

The closed-loop formulation changes the role of \(c_1\) and \(c_2\):
\begin{itemize}
  \item \(\alpha_1\) and \(\alpha_2\) define the torsion target only.
  \item \(c_1\) and \(c_2\) are free \(\phi\)-scale parameters selected by a
  reduced optimization problem.
  \item The current v2 histogram fit can be used as an initial guess, but not
  as the final definition of optimality.
  \item A fair test of the idea should compare \(J_h(c_1,c_2)\), final mass,
  active overlap, and density mismatch after the semilinear PDE has converged.
\end{itemize}

The key mathematical object is therefore the reduced map
\[
  (c_1,c_2)
  \longmapsto
  \phi_h(c_1,c_2),
\]
where \(\phi_h(c_1,c_2)\) is obtained by solving the nonlinear PDE.  Any
selection rule that does not evaluate this map is a surrogate and should be
validated against the closed-loop objective.

\end{document}
